\section{Independent review of the binary surface product extension}
\label{r20audit:PX20supported}

\paragraph{Verdict and exact version.}
\textbf{PASS relative to the explicitly frozen actual Green,
alteration, inertia, and descent interfaces.}
The reviewed source is
\texttt{inertia/product\_extension\_probe.tex}, SHA--256
\texttt{b618a9e134f13e78b19f1e8a90c7a566e2f87f6138caf93765bbf3500accee91}.
The reviewer is the supported20 branch. This is an internal independent
derivation check; the author and reviewer communicated about two
precision corrections, both present in this version. It is not external
peer review. No result for three or more bad surface factors is certified.
Historical R9 sources were used only to identify the old candidate and
its scope; none is imported here as a frozen authority.

\subsection{The chosen global projective subdivision}

The product local lattice has rank $r+s-1$, with the integral height
function $\sum a=\sum b$. Its rays $n_{ij}$ are primitive and have
height one. Along a monotone path, the first ray and the successive
differences are an integral basis: the differences supply all adjacent
first-factor differences and all adjacent second-factor differences,
and the first ray supplies the height direction. Thus every maximal
cone is unimodular and its local model is reduced strictly semistable.
No new ray, hence no new vertical component, is introduced.

The source does not use an unspecified Hartl blowup order. The heights
$-ij$ give a coherent staircase subdivision. For incomparable vertices,
the replacement by their coordinatewise ordered pair lowers the sum
of the heights by $(i'-i)(j'-j)>0$, while preserving the marginal
coordinates. Hence every lower face has chain support. Monotone
coupling of marginals gives the reverse implication, and generic
marginals realize all maximal paths. The heights are defined from
global component orders, so their restrictions agree on common faces.
The resulting integral relatively strictly convex support function
gives projectivity of this toroidal modification. This supplies the
projectivity needed later; unimodularity alone would not do so.

The reviewer independently opened Hartl's primary paper, including
Proposition~2.1 and the two monomial blowup charts on printed
pp.~218--219, and Proposition~2.2 for ramified base change. The charts
match the product cone calculation. They are not used as a substitute
for the additional component-tower calculation proved in PX20.

\subsection{Component blowup towers and genuine higher strata}

Modulo the ray $n_{ij}$, the rays from the first and second coordinate
directions become $u_a,v_b$, and the other ray becomes $u_a+v_b$.
A monotone path containing $(i,j)$ has vertices only in the lower and
upper coordinate rectangles. The lower and upper ray sets are
disjoint. On either rectangle the elementary star-subdivision
induction is valid: after inserting $u_m+v_n$, its two maximal regions
are the problems of sizes $(m,n-1)$ and $(m-1,n)$ with that final ray
adjoined. A remaining original two-face is still a face until it is
itself subdivided. Thus every subsequent center is a smooth
codimension-two stratum.

Reversing the lower coordinate orders gives the stated ascending lower
and descending upper lists. These operations globalize as the strict
transforms of the indicated products of boundary curves. Their
reduced boundary remains SNC at every stage. In particular the
intersection of the relevant two strict-transform divisors is the
strict transform of the original center, with its connected components
treated separately; it does not acquire an unlisted codimension-two
center over a higher stratum.

The component map is therefore an actual sequence of smooth
codimension-two blowups of $S_i\times T_j$. Each center maps properly
birationally to a curve product. The reviewer separately checked
Stacks Project, Tag~0C5R: a proper birational morphism between regular
integral surfaces factors into closed-point blowups, with no
characteristic-zero assumption. Consequently each center has a
curve-product motive plus point-motive corrections. The ordinary
smooth blowup formula yields
\[
 h(P_{ij})\simeq h(S_i\times T_j)\oplus K_{ij}
\]
as a Chow-motive decomposition, with actual algebraic inclusion and
projection maps and base projector $f_{ij}^*f_{ij*}$.

As an additional finite verification, the reviewer wrote
\texttt{supported/verify\_binary\_component\_fans.py}. It independently
enumerates all monotone paths and compares their component stars
with the stated two-face subdivisions for every $1\le r,s\le3$ and
every component label. It checks $36$ stars and $70$
positive-dimensional maximal cones, including all nine stars of
triple times triple, and verifies each determinant is $\pm1$.
This is a check of the local calculation, not a replacement for the
global gluing proof. At triple times triple the product has actual
higher intersections; no ordinary-node assumption has entered.

\subsection{Same-label edges and full exceptional transport}

For a double joining $P_{ij}$ and $P_{il}$, the first-factor
coordinates belonging to every other component have valuation zero
on both generating rays. They are invertible on the generic orbit
of the two-cone. Therefore no extra double component can lie wholly
over the first-factor boundary. Over the complement of that boundary
the model is the unchanged product, including at old triple points
in the second factor. This proves the single degree-one birational
map to the relevant component of $S_i\times(T_j\cap T_l)$.
An edge changing the first label maps into that label's boundary.
The symmetric statements hold for the second label.

The supported-cohomology transport lemma treats all inputs. For
$\beta=f^*f_*\beta+\beta_0$, the base summand follows from
$h^*p=p'h^*$. On the second summand, $f_*\beta_0=0$ and
$\beta_0$ vanishes off $f^{-1}Z$. Localization makes its restriction
and subsequent pushforward supported on $Z'$, where $p'$ kills it.
Thus $g^*P=P'g^*$ on full cohomology, not only on the base summand.
Poincare transposition gives the Gysin formula. It is essential here
that the support-annihilation hypothesis is imposed on $p'$ too.

For the surface projector $p_S$, its only nonzero realization is the
orthogonal complement of the geometric divisor classes in $H^2(S)$.
The degree-two supported image of an SNC curve boundary is generated
by its divisor classes; point intersections have no such contribution.
Restriction of $H^2(S)$ to that boundary is detected on its smooth
curve components and is the same divisor pairing. Therefore both
annihilations in (PX20.9) hold. Kunneth with supports carries this to
the product coefficient. No statement that all Chow maps from
$t_2(S)$ to arbitrary curve motives vanish is used.

\subsection{The four sectors are sufficient without Chow commutation}

The reviewer checked the surface projectors in Kahn--Murre--Pedrini,
Propositions~7.2.1 and~7.2.3, directly in the primary author manuscript.
They give self-dual rational Chow projectors. The geometric
Neron--Severi summand is Artin--Tate over the native field. The final
PX20 version explicitly allows a further finite unramified extension
splitting these finitely many actions and the finite point-motive
corrections. Thus the source's later description by Tate and
Picard/Albanese pieces is justified over the working residue field.
It is not inferred merely from splitting component constants.

Since $f_*f^*=1$ as correspondences, $P^S$ and $P^T$ are commuting
Chow idempotents. Their product and complementary products therefore
define genuine Chow summands. Boundary transport is required only on
cohomology: edges changing a factor label kill that factor's
transcendental projector; same-label edges have the common transported
projector on the double. This proves cohomological commutation with
the full incidence operator $A$, which is all the splitting argument
requires. It does not silently assert Chow-level commutation.

On a same-label double with map $f'$ and endpoint embeddings $h_j,h_l$,
the transported off-diagonal entry is
\[
 f_{ij*}g_{j*}g_l^*f_{il}^*
     =h_{j*}f'_*f'^*h_l^*=h_{j*}h_l^*.
\]
The middle equality uses the actual degree-one birational map.
The diagonal is the negative sum of the corresponding
$g_*g^*$ terms, so the identical calculation includes it.
Changing-label terms vanish on the selected transcendental sector.
This proves (PX20.12) with the full surface incidence operator,
including its signs and all boundary terms.

The double-transcendental sector is zero for $A$. On either
single-transcendental sector, PI20's surface inverse supplies the
inverse in the other factor; inserting its complementary projector
preserves $AGA=A$ because the surface incidence annihilates $p_S$
on both sides. The exceptional summand is killed by $P^S$ and $P^T$
as Chow projectors, and hence lies in $Q$. Its motive is already
identified by the actual blowup tower. The base of $Q$ is the tensor
product of the two nontranscendental surface parts. Consequently
$Q$ is an explicitly identified abelian-type Chow summand; it is not
an unspecified complement in a conjectural semisimple homological
category.

\subsection{One prime, numerical comparison, and actual correction}

The finitely many abelian carriers of $Q$ may be combined into one
product. Its powers, Tate twists, and finite direct sums contain all
correspondence corners needed for $QAQ$ and an opposite-degree inverse.
Deligne--Clozel, Theorem~1 and Remark~(b), were read in the primary
paper: a single suitable invertible prime works for every power of
this one fixed abelian variety. This avoids an unjustified
intersection of several unrelated positive-density prime sets.

Jannsen is applied only in numerical motives, giving a numerical
inner inverse of the corner $QAQ$. Choose actual Chow representatives
for its entries. Their error is a numerical-zero correspondence in
the same corner. Transporting it into the selected abelian carrier,
using the actual summand inclusion and projection, makes its
geometric class zero at the selected prime. Transport back kills
the entire corner error. No comparison for all correspondences on
the unrestricted fourfold component products is assumed.

The native-to-geometric numerical step is also valid. A geometric
test cycle descends to some finite constant extension. For a native
numerical-zero class, intersection with the norm of this test is
zero; projection formula identifies this with the full extension
degree times the geometric intersection. Therefore the class remains
numerically zero geometrically. All these intersections are on proper
finite-field varieties. No arithmetic-model $\operatorname{CH}_0$
degree is used.

Summing the three sector inverses yields an actual degree-minus-one
correspondence $G$ with $A^HG^HA^H=A^H$ on all cohomology.
If auxiliary representatives require a further finite residue
extension of the fixed model, normalized Galois averaging is valid
because this equation is linear in $G$. Field changes of the model
itself are handled separately by the final valued-field descent.
Only one suitable prime is asserted for the abelian core.

For a native numerical kernel vector with actual representative $x$,
the class $A^Hx$ is numerically zero. Every actual correspondence
preserves numerical equivalence, so $G^HA^Hx$ is still numerically
zero. The correction
\[
             x^\flat=x-G^HA^Hx
\]
has the original numerical class and $A^Hx^\flat=0$. This is exactly
the AK19 actual lifting requirement. Tensoring the correspondence
identity with $\Delta_{V_0}$ retains full cohomology, so arbitrary
mixed actual classes are included. No generation of their Chow groups
by external products is assumed.

\subsection{The local theorem and the uncrossed boundary}

After a common finite valued extension, the good coefficient model
survives the later extensions. Hartl's ramified base-change repair is
applied to the two surface factors before their staircase product.
All surface cohomology degrees satisfy the required
weight--monodromy input; the good factor has pure unramified
cohomology. Kunneth and the tensor rule for monodromy filtrations
give precisely OW19's odd-degree hypothesis for this witness.
Thus the actual inertia image is zero. Combining this with the
actual-lift correction makes the pulled-back cycle's periods zero.

The dominant alteration step and the finite constant-field descent
are distinct frozen interfaces. The latter uses the full extension
degree and the original extending norm. FD17/VT17 then recover the
ordinary actual Green form on the originally specified complement
with every actual compact test. No proper actual-form pushforward
and no finiteness on that complement is asserted. The final version
states $1\le r\le\dim(V\times S\times T)$ explicitly. Fresh
applications after finite valued extensions include new cycles that
need not descend.

The reviewed new proof does not enlarge the already frozen
four-dimensional result for two surfaces alone. Its good-coefficient
extension verifies the higher-dimensional mixed-cycle candidate
under the present interfaces, with honest R9 provenance. It does
not prove that a resolved product of three bad factors has the
same tower or residual core, nor does it turn its incidence map
into a naive tensor sum. General strong G and the full height
domain for arbitrary varieties remain open.
